\documentclass[aspectratio=169,10pt,spanish]{beamer}

\input{../../../_preambulo_beamer.tex}
\input{../../../_comandos_beamer.tex}

\selectlanguage{spanish}

\title{MA1001B - Modelación Estadística para la Toma de Decisiones}
\subtitle{Lección 41: Prueba ji-cuadrada de independencia}
\author{}
\institute{Tecnol\'ogico de Monterrey}
\date{}

\begin{document}

\begin{frame}[plain]
  \vspace{1.2cm}
  {\Large\bfseries MA1001B - Modelación Estadística para la Toma de Decisiones\par}
  \vspace{0.7cm}
  {\large Lección 41: Prueba ji-cuadrada de independencia\par}
  \vspace{0.7cm}
  {\color{TecRojo}\rule{\textwidth}{1.2pt}}
  \vspace{0.8cm}

  Facultad MA1001B\\[0.3cm]
  Periodo\\[0.3cm]
  Tecnol\'ogico de Monterrey
\end{frame}

\begin{frame}{La pregunta de asociación de doble entrada}
  \begin{alertblock}{Pregunta guía}
    En una muestra de $n=325$ visitantes, ¿es la conversión independiente
    del canal de adquisición?
  \end{alertblock}

  \vspace{0.6em}
  Al final de esta lección, usted puede:
  \begin{itemize}
    \item formar conteos esperados a partir de los márgenes de la tabla;
    \item calcular $\chi^2$ con $df=(r-1)(c-1)$;
    \item reportar un p-valor y la V de Cramer;
    \item rechazar una lectura causal de una asociación significativa.
  \end{itemize}

  \vfill
  \scriptsize Todos los conteos de visitantes usados hoy son sintéticos. No se usan datos reales de empresa.
\end{frame}

\begin{frame}{Una muestra, dos clasificaciones}
  \small
  \begin{center}
    \begin{tabular}{lrrrr}
      \toprule
      \textbf{Resultado} & \textbf{Correo} & \textbf{Búsqueda} & \textbf{Redes} & \textbf{Total} \\
      \midrule
      Convertido & 40 & 55 & 18 & 113 \\
      No convertido & 80 & 70 & 62 & 212 \\
      \midrule
      Total & 120 & 125 & 80 & 325 \\
      \bottomrule
    \end{tabular}
  \end{center}

  \vspace{0.4em}
  \begin{exampleblock}{Conteo esperado de independencia}
    $E_{ij}=(\text{total de fila}\times\text{total de columna})/n$. Conteo
    esperado convertido-correo: $113\times120/325=41.723$.
  \end{exampleblock}
\end{frame}

\begin{frame}[fragile]{Paso 1: Tabla de contingencia}
  \begin{columns}[T]
    \begin{column}{0.42\textwidth}
      \small
      \begin{lstlisting}
E = outer(row, col) / n
      \end{lstlisting}

      \begin{block}{Salida verificada}
        $n=325$\\
        $E$ convertidos: $41.723$, $43.462$, $27.815$\\
        Todos $E\ge 5$: sí
      \end{block}
    \end{column}
    \begin{column}{0.58\textwidth}
      \centering
      \includegraphics[width=\textwidth]{../figuras/independencia_01_tabla_contingencia.png}
      \vspace{0.2em}
      \scriptsize Tabla observada $2\times 3$ del Paso 1
    \end{column}
  \end{columns}
\end{frame}

\begin{frame}{Independencia $\chi^2$ con $df=2$}
  \small
  \[
    \chi^2=\sum\frac{(O_{ij}-E_{ij})^2}{E_{ij}}=10.114985,
    \qquad df=(2-1)(3-1)=2.
  \]
  \[
    p=0.006361,\qquad
    V=\sqrt{\chi^2/(n\min(r-1,c-1))}=0.176417.
  \]
  La prueba rechaza $H_0$: conversión independiente del canal. La V de
  Cramer es modesta.
\end{frame}

\begin{frame}[fragile]{Paso 2: Decisión ji-cuadrada}
  \begin{columns}[T]
    \begin{column}{0.40\textwidth}
      \small
      \begin{lstlisting}
chi2, p, df, E = chi2_contingency(table)
V = sqrt(chi2 / (n * (k - 1)))
      \end{lstlisting}

      \begin{block}{Salida verificada}
        $\chi^2=10.114985$, $df=2$\\
        $p=0.006361$, $V=0.176417$\\
        Decisión: rechazar $H_0$
      \end{block}
    \end{column}
    \begin{column}{0.60\textwidth}
      \centering
      \includegraphics[width=\textwidth]{../figuras/independencia_02_chi_cuadrada.png}
      \vspace{0.2em}
      \scriptsize Contribuciones celulares del Paso 2
    \end{column}
  \end{columns}
\end{frame}

\begin{frame}{Paso 3: Asociación no es una afirmación causal}
  \begin{columns}[T]
    \begin{column}{0.42\textwidth}
      \small
      $p$ conjunto $=0.006361$ (rechazar independencia).

      \vspace{0.3em}
      Estrato escritorio $p=0.368546$.

      \vspace{0.3em}
      Estrato móvil $p=0.761939$.

      \vspace{0.5em}
      \textbf{Límite:} el canal no se asignó al azar. Tras estratificar por
      dispositivo, la asociación desaparece. Un $\chi^2$ significativo no es
      una afirmación causal.
    \end{column}
    \begin{column}{0.58\textwidth}
      \centering
      \includegraphics[width=\textwidth]{../figuras/independencia_03_limite_no_causalidad.png}
      \vspace{0.2em}
      \scriptsize Confusor de dispositivo del Paso 3
    \end{column}
  \end{columns}
\end{frame}

\begin{frame}{Verificación de conceptos}
  \begin{enumerate}
    \item Calcule $df$ para una tabla $2\times 3$.
    \item ¿Por qué $E_{11}=113\times120/325$ y no $np$ de una mezcla
      hipotética?
    \item Si $p=0.006361$ y $\alpha=0.05$, ¿cuál es la decisión de
      independencia?
    \item ¿Por qué el escritorio $p=0.368546$ y el móvil $p=0.761939$
      bloquean una afirmación causal de canal?
  \end{enumerate}

  \vspace{0.6em}
  \begin{block}{Conexión con la Lección 42}
    La sesión puede continuar directamente con una prueba de homogeneidad
    ji-cuadrada, que usa el mismo estadístico con otra historia de muestreo.
  \end{block}

  \vfill
  \scriptsize Fuente: OpenStax, \textit{Introductory Business Statistics 2e},
  Capítulo 11, Sección 11.4. CC BY-NC-SA.
\end{frame}

\appendix



\end{document}
